\section{Christoffel symbols and Hamilton's equations}
\label{app:christoffel}

The nine independent non-zero Christoffel symbols of the Schwarzschild metric are
\begin{gather*}
\Gamma^t_{tr} = \frac{1}{r(r - 2)}, \qquad
\Gamma^r_{tt} = \frac{r - 2}{r^3}, \qquad
\Gamma^r_{rr} = -\frac{1}{r(r - 2)}, \qquad
\Gamma^r_{\theta\theta} = -(r - 2), \\
\Gamma^r_{\phi\phi} = -(r - 2)\sin^2\theta, \qquad
\Gamma^\theta_{r\theta} = \Gamma^\phi_{r\phi} = \frac1r, \qquad
\Gamma^\theta_{\phi\phi} = -\sin\theta\cos\theta, \qquad
\Gamma^\phi_{\theta\phi} = \cot\theta .
\end{gather*}
The code writes them in terms of $r - 2$ rather than $f = 1 - 2/r$, which keeps every
component accurate to full relative precision down to the horizon. Hamilton's equations for
(\ref{eq:hamiltonian}) are
\begin{gather*}
\dot t = -\frac{p_t}{f}, \quad \dot r = f p_r, \quad \dot\theta = \frac{p_\theta}{r^2}, \quad
\dot\phi = \frac{p_\phi}{r^2\sin^2\theta}, \qquad \dot p_t = \dot p_\phi = 0, \\
\dot p_r = -\frac{f'}{2}\left(\frac{p_t^2}{f^2} + p_r^2\right)
  + \frac{1}{r^3}\left(p_\theta^2 + \frac{p_\phi^2}{\sin^2\theta}\right), \qquad
\dot p_\theta = \frac{\cos\theta}{r^2\sin^3\theta}\,p_\phi^2 .
\end{gather*}
Both sets were checked against a symbolic derivation from $g_{\mu\nu}$ evaluated at 50 digits,
to $6\times10^{-16}$ relative per component for $r - 2$ from $10^{-8}$ to $10^4$.

\section{Gauss--Legendre methods}
\label{app:gauss}

The $s$-stage Gauss method has nodes $c_i$ at the zeros of the shifted Legendre polynomial of
degree $s$, and $a_{ij} = \int_0^{c_i}\ell_j(\tau)\,\dd\tau$, $b_j = \int_0^1\ell_j(\tau)\,\dd\tau$
with the Lagrange polynomials $\ell_j$ on the nodes. For $s = 2$,
\begin{equation*}
\begin{array}{c|cc}
\tfrac12 - \tfrac{\sqrt3}{6} & \tfrac14 & \tfrac14 - \tfrac{\sqrt3}{6} \\[2pt]
\tfrac12 + \tfrac{\sqrt3}{6} & \tfrac14 + \tfrac{\sqrt3}{6} & \tfrac14 \\[2pt]
\hline
 & \tfrac12 & \tfrac12
\end{array}
\qquad\text{and for } s = 3,\qquad
\begin{array}{c|ccc}
\tfrac12 - \tfrac{\sqrt{15}}{10} & \tfrac{5}{36} & \tfrac29 - \tfrac{\sqrt{15}}{15} & \tfrac{5}{36} - \tfrac{\sqrt{15}}{30} \\[2pt]
\tfrac12 & \tfrac{5}{36} + \tfrac{\sqrt{15}}{24} & \tfrac29 & \tfrac{5}{36} - \tfrac{\sqrt{15}}{24} \\[2pt]
\tfrac12 + \tfrac{\sqrt{15}}{10} & \tfrac{5}{36} + \tfrac{\sqrt{15}}{30} & \tfrac29 + \tfrac{\sqrt{15}}{15} & \tfrac{5}{36} \\[2pt]
\hline
 & \tfrac{5}{18} & \tfrac49 & \tfrac{5}{18}
\end{array}
\end{equation*}
GL1 is the implicit midpoint rule ($c = a = \tfrac12$, $b = 1$). One step solves
$Z_i = h\sum_j a_{ij}F(y_n + Z_j)$ by fixed-point iteration and sets
$y_{n+1} = y_n + \sum_i d_i Z_i$ with $d^\mathsf{T} = b^\mathsf{T}A^{-1}$, which avoids a final
evaluation of $F$. The start value of $Z$ is the collocation polynomial of the previous step,
evaluated at the new nodes. The iteration stops when $\|\Delta Z\|$ no longer decreases and is
within 100 ulp of $\|Z\|$ (Sec.~\ref{sec:integrators}), or after 60 iterations, which never
occurred in the experiments.

\section{Reproducibility and computing environment}
\label{app:environment}

All integrators, metrics and formulations are compiled with Numba~\cite{numba}. One compiled
driver per kind of method receives the vector field as a typed first-class function, so that
compilation happens once per machine (about 60~s) and is cached. Every run is described by a
specification (test case, formulation, method, options) whose hash keys a cache of result rows
that carry the git commit, library versions and CPU. Long runs are reduced to per-period
maxima before caching. The command \texttt{make all} regenerates every figure and table from
an empty cache and then compiles this report; on the machine of \tab{environment} it takes
about 23 minutes on four cores. Repeating it reproduced every summary table and figure byte
for byte, except the wall times, which agreed to within 25\% (median ratio 0.98).

\begin{table}[h]
\caption{\label{tab:environment}Computing environment of the results (T8).}
\centering\small
% Generated by scripts/make_tables.py from results/summary/. Do not edit.
\begin{tabular}{ll}
\toprule
Item & Value \\
\midrule
CPU & Intel(R) Core(TM) i5-7300HQ CPU @ 2.50GHz \\
logical cores & 4 \\
operating system & Linux-7.0.0-34-generic-x86\_64-with-glibc2.43 \\
Python & 3.14.7 \\
numpy & 2.5.3 \\
scipy & 1.18.1 \\
numba & 0.68.0 \\
llvmlite & 0.50.0 \\
mpmath & 1.3.0 \\
matplotlib & 3.11.2 \\
pandas & 3.0.6 \\
geoint & 1.0.0 \\
threads per process & 1 (NUMBA\_NUM\_THREADS = OMP\_NUM\_THREADS = 1) \\
parallel sweeps & up to 4 forked worker processes; timing runs one at a time \\
timing & median of 2 repeats after a warm-up run (ADR 0003) \\
floating point & IEEE 754 double; compensated summation in fixed-step updates \\
\bottomrule
\end{tabular}

\end{table}
